\section{Limitations}\label{sec:limitations}

\paragraph{Outcomes.} Most metrics are retrieval-level. The answer test rests on one small reader, short extractive questions and small samples, so its nulls do not show that methods are equal, and a stronger model may use notes differently. The chains measure a real agent on small documentation chores, not on real tasks, and their answers are scored by string match on a name or value. How much recall per token turns into finished work is still unknown.

\paragraph{Baselines.} Mem0 with inference, Graphiti, A-MEM, HippoRAG, LangMem and Cognee all call an LLM on write, so the comparison does not cover them, and it says nothing about what LLM extraction would add to those systems. Published evidence suggests a small effect on the write side \cite{yuan2026diagnosing}.

\paragraph{Generalisation.} The relevance model is fitted on four benchmark families, with BEAM-1M the only source of long conversations, so the BEAM results concern new conversations from a family the model has seen, and no untouched conversation benchmark is left. LongMemEval-S and LoCoMo informed the design and are not clean. Exact code facts are read only for Python, C and C++, so nothing here says how they behave in TypeScript, Go or Rust, and each of those languages rests on three repositories. A deployment on notes unlike these may behave differently. Beyond English, only Russian was measured, on XQuAD and not on coding notes.

\paragraph{RepoMem is a proxy.} Queries come from templates over extracted facts, so they contain identifiers and favour lexical search, while real prompts are vaguer, and on plain-language questions the gate and the weight floor hide most notes (Section~\ref{sec:verdicts}). The agent's notes are simulated from code and commit history, and every changed fact is assumed to be observed. The task track uses file overlap as relevance, which is a weak label. Keyed mode gets perfect slot keys from the static analyser, so it is an upper bound on what a real write path provides, and plain mode is the conservative end.

\paragraph{Validity.} The near-twin rule works where versions differ only in a value. It does not fire on natural-language updates, and without keys it still lets an outdated version through \V{tr-third-stale-kanerva}\% of the time on the main test. Those updates depend on a host agent, which was measured on one model and a few dozen notes, and which sometimes retired more than it should. Without vectors the list misses most paraphrases, and it appears on many notes that replace nothing, each time costing the agent tokens.

\paragraph{The gate.} It costs recall on a store's own questions, and those queries get nothing when it closes. Its false-alarm rates come from foreign questions, and on one developer's real prompts it separates far less. How often it is wrongly silent in real use is not known.

\paragraph{Learning.} A verdict covers only notes the agent was shown, so the shadow check cannot see what either ordering left out. Verdicts are tied to a project rather than a session, so two sessions open on one project can mix. The signal test labels notes by construction, and its distractors are easy to reject. The chains use one model, small chores and six repositories at a time. A project's gate moved in only two of them, late in the chain, so per-project tuning over longer use is unmeasured.

\paragraph{Plugins.} The benchmarks ran the library with exact token counts, while the plugins estimate tokens with a dependency-free approximation. Until the embedding model arrives, a new install runs the lexical model, so its first sessions get the lexical numbers.
